\chapter{All tests for the uniform infinitesimal C1 source}
\label{component:p6cad7544d1b8}
\noindent\textbf{Component }\path{p6cad7544d1b8}\par
{\small\noindent Source: \path{CORE/prior/inherited/c1_uniform_all_tests.tex}\par}
{\small\noindent SHA256: \path{6cad7544d1b8537e36bef0527c8f2a6dd53113d9a6ba10856bea98dc9142a113}\par}
{\small\noindent Inventory use: Two strong currents and Critical current mass lemmas, before the uniform theorem | C1\_\allowbreak{}BARE\_\allowbreak{}OTHER\_\allowbreak{}TERMINALS; precise subordinate uses retained below\par}
\medskip
\section{All tests for the uniform infinitesimal C1 source}
\label{p6cad7544d1b8:sec:c1-uniform-all-tests}

The two state charts can be joined without importing an all-tests atlas
label. We first record two short outer-angle identities.

\begin{lemma}[Two strong currents]
Suppose the original event coefficients are nonnegative, $d>0$, and
$d/c,d/D$ are infinitesimal, where $c=e_3$, $D=e_5$.
Then every testing chart with $1/v$ bounded is tensor-SOS, for every
other angle and every state.
\end{lemma}
\begin{proof}
Put $g=cD/(c+D)$, $h=D/(c+D)$, $R=uW-Y$,
$C=v(uX-Z)$, $B=X+uZ$, and $N=R+hC$.
Let $D_u=1+u^2$, $D_v=1+v^2$ and set
\[
 K_0=2\sqrt g\,C,\qquad
 B_0=\sqrt{2d}\sqrt{D_v/D_u}\,B,\quad
 c_u=D_u^{-1/2},\quad t_u=uD_u^{-1/2},\quad
 \eta_v=\sqrt{D_v}/v.
\]
The complete selected parent is exactly
\[
\begin{split}
4cD_vR^2+4D(R+C)^2+d\mathcal H
={}&4(c+D)N^2+4cv^2R^2+2d(R+C)^2\\
 &+2dD_u(vY-Z)^2+(K_0,B_0)Q(K_0,B_0)^T,
\end{split}
\]
where
\[
 Q=\begin{pmatrix}
 1-\dfrac{d}{2g}c_u^2\eta_v^2&
 \sqrt{\dfrac{d}{2g}}t_u\eta_v\\
 \sqrt{\dfrac{d}{2g}}t_u\eta_v&1
 \end{pmatrix}.
\]
To check the identity, use $Z=(uB-C/v)/D_u$ and complete the two current
squares. Since $d/g=d/c+d/D$ is infinitesimal and $c_u,t_u,\eta_v$
are bounded, $\bar Q=I_2$. Restore every unselected positive event.
All roots separate between fields. Clearing the original pure-testing
homogeneous angle factors gives the same identity at literal infinity:
$\eta_v^2=1+(1/v)^2$, and $c_u,t_u$ are the normalized homogeneous
Alice coordinates. Thus no infinite-angle endpoint is excluded.
\end{proof}

\begin{lemma}[Critical current mass outside the small Alice chart]
Under the preceding strength assumptions, suppose
\[
 \gamma=\frac{A g}{d^2}\quad\hbox{is bounded},\qquad \bar\gamma\geq\tfrac14,
\]
where $A=e_4$ is the original event coefficient. Then every independent
test with $v$ infinitesimal and $1/u$ bounded is tensor-SOS, for every
state, including the homogeneous Alice-infinite endpoint.
\end{lemma}
\begin{proof}
Keep $R,C,B,N,g,h,c_u,t_u$ from the preceding proof, and put
$\epsilon=d/g$, $q=\sqrt\epsilon$ and
\[
 K_1=2\sqrt g\,C+\sqrt{d\epsilon}\,Y,\quad
 B_1=\sqrt d(B-uvY),\quad V_0=\sqrt d\,vY,
 \quad Y_0=\sqrt{d\epsilon}\sqrt{D_u}\,Y.
\]
Define
\[
 L_0=\frac{qc_u}{2}(K_1-c_uY_0)-t_uvB_1+c_uvV_0.
\]
Direct multiplication gives $L_0=\sqrt{dD_u}\,v(vY-Z)$ and the exact
full selected-parent identity
\[
\begin{split}
&4cD_vR^2+4D(R+C)^2+4AD_uD_vY^2+d\mathcal H\\
&\qquad=4(c+D)N^2+4cv^2R^2+2d(R+C)^2+\mathcal G,\\
\mathcal G={}&K_1^2+2D_v(B_1^2+V_0^2)
 +(4\gamma D_v-c_u^2)Y_0^2\\
&+2qv(K_1-c_uY_0)V_0-2L_0^2.
\end{split}
\]
The four-root Gram of $\mathcal G$ has residue
$\operatorname{diag}(1,2,2,4\bar\gamma-\bar c_u^2)$.
Its last entry is at least $1-\bar c_u^2=\bar t_u^2>0$ because $1/u$
is bounded. Every correction is polynomial in bounded
$q,v,c_u,t_u,\gamma$ and has zero residue. The strict-residue lemma and
restoration of the other events prove the claim. The displayed bounded
homogeneous factors also cover literal Alice infinity.
\end{proof}

\begin{theorem}[Whole uniform infinitesimal C1 domain]
\label{p6cad7544d1b8:thm:c1-uniform-all-tests}
Use the unchanged canonical uniform C1 source, including all exact
coquotients and original coefficient zeros. Suppose
\[
 k,f,g,e>0\text{ are infinitesimal},\qquad e=kf/g,
\]
and the positive bounded shapes satisfy
\[
 R=1+k^2(1-af^2),\quad T=1+k^2-k^2g^2t,\qquad
 M=R-t>0,\quad J=T-a>0.
\]
Assume every original event coefficient is nonnegative. Then the entire
parent is tensor-SOS at every independent projective measurement pair
and every real or complex testing state. There is no restriction on
$b=g/k$ or on $A_c=g/f$.
\end{theorem}

\begin{proof}
First derive the two outer-angle premises from the original coefficients.
Write $B_c=2c_5-Jf^2c_d$. The exact source gives
\[
 D=\frac{aB_c}{Jc_3},\qquad
 \frac dD=\frac{2Jf^2c_d}{B_c},\qquad
 \frac dc=\frac{af^2c_d}{c_3}.
\]
Since $\bar B_c=4$, both ratios are infinitesimal, independently of
any loss of residue of $a$ or $J$. Also, for the original harmonic
mass $g_B=cD/(c+D)$,
\[
 \frac{A g_B}{d^2}
 =\frac{c_4c_3B_c}{c_d^2(2Jc_3+aB_c)}.              \tag{UA1}
\]
The denominator is a unit: its residue is
$16\cdot4(\bar J+\bar a)=64$, since $J+a=T$ has residue one.
The numerator has residue $16$. Thus the original mass in (UA1) has
residue exactly $1/4$. This argument uses neither a strong first pair
nor beta positivity beyond the original source hypotheses.

The two preceding lemmas handle every test except $u,v$ infinitesimal.
Literal $u=0$ or $v=0$ also gives a nonnegative pure-testing core
immediately. Normalize a nonzero state by a coordinate of largest
absolute value; all coordinates are then bounded and at least one is
a unit. If $X$ or $Y$ is a unit, the pure-testing core has positive
residue $2(X^2+Y^2)$, so its product with the positive coefficient $d$
is itself a square and all other events are positive squares.
It remains that $X,Y$ are infinitesimal.

The following coefficient-only split is exhaustive. If $b=g/k$ is
bounded, then $f=eb$ and $g=kb$, so the bounded-$b$ Z theorem treats
$Z$ a unit and arbitrary bounded $W$. If $Z$ is infinitesimal, $W$ is
a unit, and the bounded-$b$ all-rho W endpoint applies after $W=1$
normalization.

If $b$ is unbounded, put $A=f$, $H=g$, $K=k/g=1/b$.
All three are infinitesimal. Original beta nonnegativity gives exactly
\[
 \beta+\frac d2=\frac{2ae^2c_b}{tc_3},\qquad
 \frac d2=\frac{af^2c_d}{c_3}
 \quad\Longrightarrow\quad
 E:=tb^2\leq\frac{2c_b}{c_d}.
\]
Hence $E>0$ is bounded with $\bar E\leq1$ and the exact source map is
\[
 (f,g,k,e,a,t)=(A,H,HK,AK,a,K^2E).
\]
Both original gaps, every coefficient and every literal coefficient
zero are preserved. The new inverse-$b$ Z theorem handles $Z$ a unit;
the finite ten-root W theorem handles the complementary $W$-unit,
$Z$-infinitesimal chart. Thus all state cases are exhausted.

The zero state is trivial. For a complex state the real parent matrix
splits its evaluation into the sum of its real and imaginary quadratic
evaluations. All state normalizations and homogeneous projective factors
are pure testing elements. No comparison of an independent testing
quantity with a coefficient scale was used.
\end{proof}
